\clearpage
\section{대각합 쌍, 쌍대기저와 분리성}
\label{sec:trace-pairing-separability}

\begin{learninggoal}
대각합으로부터 대칭 쌍선형형을 만들고, 그 비퇴화성이 체확장의 분리성과 정확히 일치함을 증명한다. 쌍대기저와 기저 판별식을 계산하며 Chapter 2의 판별식 이론과 연결한다.
\end{learninggoal}

유한확장 $L/K$에 대해
\[
  \langle x,y\rangle_{L/K}
  :=\Tr_{L/K}(xy)
\]
로 두자. 대각합의 선형성 때문에 이는 $K$-쌍선형이고,
체의 곱셈이 가환이므로 대칭이다.

\begin{definition}[대각합 쌍]
\label{def:trace-pairing}
쌍선형형
\[
  L\times L\longrightarrow K,
  \qquad
  (x,y)\longmapsto\Tr_{L/K}(xy)
\]
를 $L/K$의 \emph{대각합 쌍}(trace pairing)이라 한다.
\end{definition}

\begin{theorem}[분리성의 대각합 판정]
\label{thm:trace-pairing-separability}
유한확장 $L/K$에 대해 다음은 동치이다.
\begin{enumerate}[label=(\roman*)]
  \item $L/K$는 분리확장이다.
  \item $\Tr_{L/K}:L\to K$는 영사상이 아니다.
  \item 대각합 쌍은 비퇴화이다.
  \item $L$의 어떤 $K$-기저 $x_1,\dots,x_n$에 대해
  \[
    \det\bigl(\Tr_{L/K}(x_ix_j)\bigr)\ne0.
  \]
  \item $L$의 모든 $K$-기저에 대해 위 행렬식이 영이 아니다.
\end{enumerate}
\end{theorem}

\begin{proofidea}
비분리확장에서는 대각합 자체가 영이다. 분리확장에서는
\[
  \Tr_{L/K}=\sigma_1+\cdots+\sigma_n
\]
이고 서로 다른 체매장의 선형독립 때문에 이 합은 영사상이 아니다.
영이 아닌 $x$가 대각합 쌍의 왼쪽 핵에 들어가면 $xL=L$이므로
대각합 전체가 영이 되는 모순이 생긴다.
\end{proofidea}

\begin{proof}
(i)$\Rightarrow$(ii)를 보자. 분리확장이므로
\[
  \Tr_{L/K}=\sigma_1+\cdots+\sigma_n
\]
이다. 서로 다른 체매장은 따름정리
\ref{cor:field-homomorphism-independence}에 의해 선형독립이므로
이 합은 영사상이 아니다.

(ii)$\Rightarrow$(iii)를 보자. 어떤 $x\in L$가
\[
  \Tr_{L/K}(xy)=0\qquad(y\in L)
\]
을 만족한다고 하자. $x\ne0$이면 $xL=L$이므로
$\Tr_{L/K}$가 $L$ 전체에서 영이 되어 (ii)에 모순이다.
따라서 $x=0$이고 쌍은 비퇴화이다.

(iii), (iv), (v)의 동치는 쌍선형형의 그람행렬에 대한 표준 선형대수이다.
기저변환행렬이 $C$이면 그람행렬은
\[
  G\longmapsto C^{\mathsf T}GC
\]
로 바뀌므로 행렬식은 $\det(C)^2$배만큼 변한다.

마지막으로 $L/K$가 비분리확장이면
따름정리 \ref{cor:inseparable-trace-zero}에서 대각합이 영이므로
(ii)--(v)는 모두 성립하지 않는다. 따라서 동치가 완성된다.
\end{proof}

\begin{corollary}[대각합의 상과 핵]
\label{cor:trace-kernel-dimension}
$n=[L:K]$라 하자.
\begin{enumerate}[label=(\roman*)]
  \item $L/K$가 분리이면 $\Tr_{L/K}$는 전사이고
  \[
    \dim_K\ker\Tr_{L/K}=n-1.
  \]
  \item $L/K$가 비분리이면
  \[
    \ker\Tr_{L/K}=L.
  \]
\end{enumerate}
\end{corollary}

\begin{proof}
분리인 경우 대각합은 영이 아닌 $K$-선형함수이고 공역 $K$는
$1$차원이므로 전사이다. 차원정리를 적용한다. 비분리인 경우는
따름정리 \ref{cor:inseparable-trace-zero}이다.
\end{proof}

\subsection{쌍대기저}

\begin{corollary}[대각합 쌍대기저]
\label{cor:trace-dual-basis}
$L/K$가 유한 분리확장이고 $x_1,\dots,x_n$이 $K$-기저라 하자.
그러면 유일한 $K$-기저 $y_1,\dots,y_n$이 존재하여
\[
  \boxed{
  \Tr_{L/K}(x_i y_j)=\delta_{ij}}
  \qquad(1\le i,j\le n)
\]
을 만족한다. 이를 $(x_i)$의 \emph{대각합 쌍대기저}라 한다.
\end{corollary}

\begin{proof}
비퇴화 쌍은 사상
\[
  L\longrightarrow L^\vee,
  \qquad
  x\longmapsto\bigl(y\mapsto\Tr_{L/K}(xy)\bigr)
\]
을 동형으로 만든다. $L^\vee$에서 $(x_i)$에 대한 표준 쌍대기저를
역상으로 옮기면 원하는 $y_j$를 얻고 유일성도 따른다.
\end{proof}

\begin{example}[이차확장의 쌍대기저]
$\charac K\ne2$, $L=K(\sqrt d)$라 하자. 기저
\[
  x_1=1,\qquad x_2=\sqrt d
\]
의 대각합 그람행렬은
\[
  G=
  \begin{pmatrix}
    \Tr(1)&\Tr(\sqrt d)\\
    \Tr(\sqrt d)&\Tr(d)
  \end{pmatrix}
  =
  \begin{pmatrix}
    2&0\\
    0&2d
  \end{pmatrix}.
\]
따라서 쌍대기저는
\[
  \boxed{
  \frac12,\qquad
  \frac{\sqrt d}{2d}}
\]
이고 기저의 판별식은
\[
  \det G=4d.
\]
\end{example}

\subsection{판별식과의 연결}

$K$-기저 $\mathbf x=(x_1,\dots,x_n)$의 판별식을
\[
  D_{L/K}(\mathbf x)
  :=\det\bigl(\Tr_{L/K}(x_ix_j)\bigr)
\]
로 쓴다. 정리 \ref{thm:trace-pairing-separability}에서
\[
  L/K\text{가 분리}
  \quad\Longleftrightarrow\quad
  D_{L/K}(\mathbf x)\ne0
\]
이다.

특히 $L=K(\alpha)$가 분리 단순확장이고
$\mathbf x=(1,\alpha,\dots,\alpha^{n-1})$이면 Chapter 2의
정리 \ref{thm:power-basis-discriminant}에 의해
\[
  D_{L/K}(1,\alpha,\dots,\alpha^{n-1})
  =\Disc(m_{\alpha,K}).
\]
따라서 최소다항식의 판별식이 영이 아닌 것, 거듭제곱기저의
대각합 행렬이 가역인 것, $\alpha$가 분리원소인 것은 같은 현상이다.

\begin{keyidea}
분리성은 ``서로 다른 켤레가 충분히 많다''는 조건이고,
대각합 쌍의 비퇴화성은 그 켤레들이 선형대수적으로 서로를 구별할 수 있다는 조건이다.
\[
  \boxed{
  \text{분리확장}
  \Longleftrightarrow
  \text{대각합 쌍 비퇴화}
  \Longleftrightarrow
  \text{기저 판별식}\ne0.}
\]
\end{keyidea}

\begin{checkpoint}
\begin{enumerate}[label=(\alph*)]
  \item $\Q(\sqrt5)/\Q$의 기저
  $(1,\sqrt5)$에 대한 대각합 행렬과 판별식을 구하라.
  \item 같은 체의 기저
  $\left(1,\frac{1+\sqrt5}{2}\right)$의 판별식을 구하라.
  \item 비분리확장에서 쌍대기저가 존재할 수 없는 이유를 설명하라.
\end{enumerate}
\end{checkpoint}
